\begin{theorem}[端点缺陷的 Poisson--Littlewood--Paley 能量律]\label{thm:xi-endpoint-defect-h12-energy-poisson-l2-dissipation}
在定理 \ref{thm:xi-cayley-modulus-poisson-fourier-fingerprint-layered} 与推论 \ref{cor:xi-cayley-depth-scan-area-lowfreq} 的有限缺陷记号下，
令
\[
\nu=\sum_{j=1}^{\kappa}m_j\,\delta_{(\gamma_j,\delta_j)},\qquad
m_j\in\NN,\ \gamma_j\in\RR,\ \delta_j\in(0,1),
\qquad
J_\nu(x):=\sum_{j=1}^{\kappa}m_j\,u_{\delta_j}(x-\gamma_j),
\]
并记总质量 \(M:=\sum_{j=1}^{\kappa}m_j\)。
定义齐次 Sobolev 半范数
\[
\|f\|_{\dot H^{1/2}(\RR)}^2:=\frac{1}{2\pi}\int_{\RR}|\xi|\,|\widehat f(\xi)|^2\,\dd\xi,
\]
其中 \(\widehat f(\xi)=\int_{\RR}e^{-\mathrm{i}\xi x}f(x)\dd x\)。
则 \(J_\nu\in\dot H^{1/2}(\RR)\)，并且成立严格恒等式
\begin{equation}\label{eq:xi-endpoint-defect-h12-energy-poisson-l2-dissipation}
\boxed{\;
\|J_\nu\|_{\dot H^{1/2}(\RR)}^2
=2\int_{0}^{\infty}\bigl\|P_t*(|D|J_\nu)\bigr\|_{L^2(\RR)}^2\,\dd t\;}
\end{equation}
其中 \(|D|\) 为 Fourier 乘子 \(|\xi|\)（命题 \ref{prop:xi-jensen-poisson-intertwiner-two-scale}）。
\par
进一步，引入两族概率密度（\(t\ge0\)）
\[
p_t(x):=\frac{1}{M}\sum_{j=1}^{\kappa}m_j\,P_{t+1-\delta_j}(x-\gamma_j),
\qquad
q_t(x):=\frac{1}{M}\sum_{j=1}^{\kappa}m_j\,P_{t+1+\delta_j}(x-\gamma_j),
\]
则
\begin{equation}\label{eq:xi-endpoint-defect-h12-energy-poisson-l2-dissipation-ptqt}
\boxed{\;
\|J_\nu\|_{\dot H^{1/2}(\RR)}^2
=2\pi^2M^2\int_0^\infty\|p_t-q_t\|_{L^2(\RR)}^2\,\dd t\; }.
\end{equation}
\leanverified{paper\_xi\_endpoint\_defect\_h12\_energy\_poisson\_l2\_dissipation}
\end{theorem}
\begin{proof}
由 Poisson 半群的 Fourier 乘子 \(\widehat{P_t*f}=e^{-t|\xi|}\widehat f\)，有
\[
\int_0^\infty\|P_t*(|D|J_\nu)\|_2^2\,\dd t
=\frac{1}{2\pi}\int_{\RR}\bigl||\xi|\widehat{J_\nu}(\xi)\bigr|^2
\left(\int_0^\infty e^{-2t|\xi|}\,\dd t\right)\dd\xi
=\frac{1}{4\pi}\int_{\RR}|\xi|\,|\widehat{J_\nu}(\xi)|^2\,\dd\xi,
\]
从而得到 \eqref{eq:xi-endpoint-defect-h12-energy-poisson-l2-dissipation}。
\par
再由命题 \ref{prop:xi-jensen-poisson-intertwiner-two-scale} 与平移线性叠加，
\[
P_t*(|D|J_\nu)
=\pi\sum_{j=1}^{\kappa}m_j\Bigl(P_{t+1-\delta_j}-P_{t+1+\delta_j}\Bigr)(\cdot-\gamma_j)
=\pi M\,(p_t-q_t),
\]
对 \(L^2\) 范数平方并积分即得 \eqref{eq:xi-endpoint-defect-h12-energy-poisson-l2-dissipation-ptqt}。证毕。
\end{proof}

\begin{proposition}[Poisson--Laplace 端点能量的任意阶多极衰减]\label{prop:xi-poisson-laplace-energy-multipole-decay}
设函数 \(\widehat D:\RR\to\CC\) 满足 \(\int_{\RR}|\xi|\,|\widehat D(\xi)|^2\,\dd\xi<\infty\)。
并假设存在最小整数 \(r\ge0\) 与常数 \(c_r\ne0\)，使得当 \(\xi\to0\) 时有
\[
\widehat D(\xi)=c_r\xi^r+o(|\xi|^r).
\]
定义 Poisson--Laplace 加权能量泛函
\[
\mathcal{E}(t):=\frac{1}{2\pi}\int_{\RR}|\xi|\,|\widehat D(\xi)|^2\,e^{-2t|\xi|}\,\dd\xi,\qquad t>0.
\]
则当 \(t\to\infty\) 时存在精确首项渐近
\[
\boxed{\;
t^{2r+2}\mathcal{E}(t)\ \longrightarrow\
\frac{|c_r|^2\,\Gamma(2r+2)}{\pi\,2^{2r+2}}\; } ,
\]
从而 \(\mathcal{E}(t)=\Theta(t^{-2r-2})\)，且首项常数完全由 \(\widehat D\) 在 \(\xi=0\) 的首个非零喷射系数 \(c_r\) 决定。
\end{proposition}
\begin{proof}
由被积函数的偶对称性，
\[
\mathcal{E}(t)=\frac{1}{\pi}\int_{0}^{\infty}\xi\,|\widehat D(\xi)|^2\,e^{-2t\xi}\,\dd\xi.
\]
作代换 \(u=t\xi\)，得
\[
t^{2r+2}\mathcal{E}(t)=\frac{1}{\pi}\int_{0}^{\infty}u^{2r+1}
\Bigl|\frac{\widehat D(u/t)}{(u/t)^r}\Bigr|^2e^{-2u}\,\dd u.
\]
由假设 \(\widehat D(\eta)/\eta^r\to c_r\)（\(\eta\to0\)）以及指数核的可积主导，可用支配收敛定理把极限移入积分，得到
\[
\lim_{t\to\infty}t^{2r+2}\mathcal{E}(t)
=\frac{|c_r|^2}{\pi}\int_{0}^{\infty}u^{2r+1}e^{-2u}\,\dd u
=\frac{|c_r|^2}{\pi}\cdot\frac{\Gamma(2r+2)}{2^{2r+2}}.
\]
证毕。
\end{proof}

\begin{theorem}[端点缺陷能量对加权 Poisson 相对熵的显式下界]\label{thm:xi-endpoint-defect-energy-kl-bridge-explicit}
\leanverified{paper\_xi\_endpoint\_defect\_energy\_kl\_bridge\_explicit}
在定理 \ref{thm:xi-endpoint-defect-h12-energy-poisson-l2-dissipation} 的设定下，
令 \(\delta_\ast:=\max_{1\le j\le\kappa}\delta_j\in(0,1)\)。
则对每个 \(t\ge0\)、\(x\in\RR\) 有一致夹逼
\begin{equation}\label{eq:xi-endpoint-defect-ptqt-ratio-bracket}
\frac{t+1-\delta_\ast}{t+1+\delta_\ast}\ \le\ \frac{p_t(x)}{q_t(x)}\ \le\ \frac{t+1+\delta_\ast}{t+1-\delta_\ast}.
\end{equation}
从而成立显式熵桥不等式
\begin{equation}\label{eq:xi-endpoint-defect-energy-kl-bridge-explicit}
\boxed{\;
\int_0^\infty \frac{1}{t+1}\,D_{\mathrm{KL}}(p_t\mid q_t)\,\dd t
\ \ge\
\frac{1-\delta_\ast}{4\pi(1+\delta_\ast)\,M^2}\,
\|J_\nu\|_{\dot H^{1/2}(\RR)}^2\; }.
\end{equation}
\end{theorem}
\begin{proof}
先证 \eqref{eq:xi-endpoint-defect-ptqt-ratio-bracket}。
对每个 \(j\)，令
\[
r_{j,t}(x):=\frac{P_{t+1-\delta_j}(x-\gamma_j)}{P_{t+1+\delta_j}(x-\gamma_j)}
=\frac{t+1-\delta_j}{t+1+\delta_j}\cdot
\frac{(x-\gamma_j)^2+(t+1+\delta_j)^2}{(x-\gamma_j)^2+(t+1-\delta_j)^2}.
\]
由 \((x-\gamma_j)^2\ge0\) 得
\[
\frac{t+1-\delta_j}{t+1+\delta_j}\le r_{j,t}(x)\le \frac{t+1+\delta_j}{t+1-\delta_j},
\]
取 \(\delta_\ast=\max\delta_j\) 并注意 \(p_t/q_t\) 是 \(\{r_{j,t}\}\) 的正权平均，即得 \eqref{eq:xi-endpoint-defect-ptqt-ratio-bracket}。
\par
令 \(u_t(x):=\frac{p_t(x)}{q_t(x)}\)。
取 \(\phi(u):=u\log u\)，则 \(\phi''(u)=1/u\ge 1/\beta_t\) 在区间 \([\alpha_t,\beta_t]\) 上成立，其中
\[
\alpha_t:=\frac{t+1-\delta_\ast}{t+1+\delta_\ast},\qquad
\beta_t:=\frac{t+1+\delta_\ast}{t+1-\delta_\ast}.
\]
因此对 \(u\in[\alpha_t,\beta_t]\) 有强凸下界
\[
\phi(u)\ge \phi(1)+\phi'(1)(u-1)+\frac{1}{2\beta_t}(u-1)^2.
\]
将 \(u=u_t(x)\) 代入并对 \(q_t(x)\dd x\) 积分，利用 \(\int q_t(u_t-1)=\int(p_t-q_t)=0\) 消去一次项，得到
\[
D_{\mathrm{KL}}(p_t\mid q_t)=\int_{\RR}q_t\,\phi(u_t)\,\dd x
\ge \frac{1}{2\beta_t}\int_{\RR}q_t(x)\,(u_t(x)-1)^2\,\dd x.
\]
另一方面，
\[
\|p_t-q_t\|_2^2=\int_{\RR}q_t(x)^2\,(u_t(x)-1)^2\,\dd x
\le \|q_t\|_{\infty}\int_{\RR}q_t(x)\,(u_t(x)-1)^2\,\dd x
\le 2\beta_t\,\|q_t\|_\infty\,D_{\mathrm{KL}}(p_t\mid q_t).
\]
又因 \(\|P_s\|_\infty=\frac{1}{\pi s}\)，故
\[
\|q_t\|_\infty
\le \frac{1}{M}\sum_{j=1}^{\kappa}m_j\,\|P_{t+1+\delta_j}\|_\infty
\le \frac{1}{\pi(t+1)}.
\]
从而
\[
\|p_t-q_t\|_2^2\le \frac{2\beta_t}{\pi(t+1)}\,D_{\mathrm{KL}}(p_t\mid q_t).
\]
两边乘以 \(\pi/(2\beta_t)\) 并对 \(t\ge0\) 积分，再用
\(
\beta_t\le\frac{1+\delta_\ast}{1-\delta_\ast}
\)
得到
\[
\int_0^\infty \frac{1}{t+1}D_{\mathrm{KL}}(p_t\mid q_t)\,\dd t
\ge \frac{\pi}{2}\cdot\frac{1-\delta_\ast}{1+\delta_\ast}\int_0^\infty \|p_t-q_t\|_2^2\,\dd t.
\]
最后由 \eqref{eq:xi-endpoint-defect-h12-energy-poisson-l2-dissipation-ptqt} 代入即得 \eqref{eq:xi-endpoint-defect-energy-kl-bridge-explicit}。证毕。
\end{proof}

\begin{corollary}[单缺陷的加权熵积分闭式与 \texorpdfstring{$\dot H^{1/2}$}{H1/2} 能量]\label{cor:xi-endpoint-single-defect-weighted-kl-li2}
在定理 \ref{thm:xi-endpoint-defect-h12-energy-poisson-l2-dissipation} 中取 \(\kappa=1\)、\(m_1=1\)、\(\gamma_1=0\)，并记 \(\delta=\delta_1\in(0,1)\)。
则
\[
p_t(x)=P_{t+1-\delta}(x),\qquad q_t(x)=P_{t+1+\delta}(x).
\]
并且对每个 \(t\ge0\) 有闭式
\begin{equation}\label{eq:xi-endpoint-single-defect-kl-closed-form}
D_{\mathrm{KL}}(p_t\mid q_t)=\log\frac{(t+1)^2}{(t+1)^2-\delta^2},
\end{equation}
以及严格积分恒等式
\begin{equation}\label{eq:xi-endpoint-single-defect-weighted-kl-li2}
\boxed{\;
\int_0^\infty \frac{1}{t+1}\,D_{\mathrm{KL}}(p_t\mid q_t)\,\dd t
=\frac12\,\operatorname{Li}_2(\delta^2)\; }.
\end{equation}
此外，端点缺陷势
\(
J_\nu(x)=u_\delta(x)=\frac12\log\frac{x^2+(1+\delta)^2}{x^2+(1-\delta)^2}
\)
满足
\begin{equation}\label{eq:xi-endpoint-single-defect-h12-energy-closed}
\boxed{\;
\|u_\delta\|_{\dot H^{1/2}(\RR)}^2=\pi\log\frac{1}{1-\delta^2}\; }.
\end{equation}
\end{corollary}
\begin{proof}
式 \eqref{eq:xi-endpoint-single-defect-kl-closed-form} 为同位置 Cauchy 分布间 KL 的闭式化简：
将
\(
p_t(x)=\pi^{-1}\frac{t+1-\delta}{x^2+(t+1-\delta)^2}
\)
与
\(
q_t(x)=\pi^{-1}\frac{t+1+\delta}{x^2+(t+1+\delta)^2}
\)
代入 \(D_{\mathrm{KL}}(p_t\mid q_t)=\int p_t\log(p_t/q_t)\) 并化简即可得到该闭式。
\par
为证 \eqref{eq:xi-endpoint-single-defect-weighted-kl-li2}，令 \(s=t+1\)，则
\[
\int_0^\infty \frac{1}{t+1}\log\frac{(t+1)^2}{(t+1)^2-\delta^2}\,\dd t
=\int_{1}^{\infty}\frac{-\log\!\left(1-\frac{\delta^2}{s^2}\right)}{s}\,\dd s.
\]
再令 \(v=\delta^2/s^2\)（即 \(s=\delta/\sqrt v\)），得
\[
\int_{1}^{\infty}\frac{-\log\!\left(1-\frac{\delta^2}{s^2}\right)}{s}\,\dd s
=\frac12\int_{0}^{\delta^2}\frac{-\log(1-v)}{v}\,\dd v
=\frac12\,\operatorname{Li}_2(\delta^2),
\]
其中末步使用 dilogarithm 的标准积分表征
\(
\operatorname{Li}_2(z)=\int_0^{z}\frac{-\log(1-v)}{v}\,\dd v
\)（\(z\in[0,1)\)）。
\par
最后由定理 \ref{thm:xi-finite-defect-poisson-l2-energy-closed-form} 的单缺陷闭式（推论 \ref{cor:xi-finite-defect-poisson-l2-energy-single-defect}，取 \(m=1\)）与
\eqref{eq:xi-endpoint-defect-h12-energy-poisson-l2-dissipation}，
\[
\|u_\delta\|_{\dot H^{1/2}}^2
=2\int_0^\infty \bigl\|P_t*(|D|u_\delta)\bigr\|_2^2\,\dd t
=2\int_0^\infty \pi\,\frac{\delta^2}{(t+1)\bigl((t+1)^2-\delta^2\bigr)}\,\dd t.
\]
令 \(s=t+1\)，并注意
\(
\frac{\delta^2}{s(s^2-\delta^2)}=\frac{s}{s^2-\delta^2}-\frac{1}{s}
\)，积分得到
\[
2\pi\int_{1}^{\infty}\left(\frac{s}{s^2-\delta^2}-\frac{1}{s}\right)\dd s
=2\pi\left[\,\frac12\log(s^2-\delta^2)-\log s\,\right]_{s=1}^{\infty}
=\pi\log\frac{1}{1-\delta^2},
\]
即为 \eqref{eq:xi-endpoint-single-defect-h12-energy-closed}。证毕。
\end{proof}

\begin{conjecture}[多缺陷加权熵泛函的 dilogarithm 交比核表述]\label{conj:xi-endpoint-multidefect-dilogarithm-crossratio-kernel}
在定理 \ref{thm:xi-endpoint-defect-energy-kl-bridge-explicit} 的设定下，定义加权熵泛函
\[
\mathfrak E_\nu:=\int_0^\infty \frac{1}{t+1}\,D_{\mathrm{KL}}(p_t\mid q_t)\,\dd t.
\]
并令上半平面点
\[
z_j^\pm:=\gamma_j+\mathrm{i}(1\pm\delta_j)\in\mathbb H\qquad(1\le j\le\kappa).
\]
猜想存在一个仅依赖交比的显式核 \(\mathcal K_{\operatorname{Li}_2}(j,k)\)
（可由 Rogers/Bloch--Wigner 型 dilogarithm 组合给出），使得
\[
\mathfrak E_\nu=\sum_{j,k=1}^{\kappa}m_jm_k\,\mathcal K_{\operatorname{Li}_2}(j,k),
\]
并满足：
\begin{enumerate}
  \item \(\mathcal K_{\operatorname{Li}_2}\) 对 \(PSL(2,\RR)\) 的 Möbius 作用不变；
  \item \(\mathcal K_{\operatorname{Li}_2}\) 为严格正定核，从而 \(\mathfrak E_\nu>0\) 当且仅当缺陷非平凡；
  \item 对角项满足 \(\mathcal K_{\operatorname{Li}_2}(j,j)=\frac12\,\operatorname{Li}_2(\delta_j^2)\)，与推论 \ref{cor:xi-endpoint-single-defect-weighted-kl-li2} 一致；
  \item 在小缺陷极限下，其二阶变分退化为对数交比核，从而与 \( \dot H^{1/2} \) Gram 核（命题 \ref{prop:xi-endpoint-jensen-defect-h12-gram-kernel}）兼容，并与 \eqref{eq:xi-endpoint-defect-energy-kl-bridge-explicit} 的桥接常数相一致。
\end{enumerate}
\end{conjecture}

\endinput
